% Cover letter for the ITOR special issue submission. Same house style
% as main.tex: no em or en dashes anywhere.
\documentclass[11pt]{article}
\usepackage[a4paper,margin=2.8cm]{geometry}
\usepackage[T1]{fontenc}
\usepackage[utf8]{inputenc}
\pagestyle{empty}
\setlength{\parskip}{0.6em}
\setlength{\parindent}{0pt}

\begin{document}

\begin{flushright}
Bogot\'a, Colombia\\
20 September 2026
\end{flushright}

\textbf{International Transactions in Operational Research}\\
Special issue on Operations Research in Wildfire Management\\
Guest editors: Prof.\ Filipe Alvelos, Prof.\ Jaime Carrasco Barra,
Prof.\ Miguel Constantino, and Prof.\ Yu Wei

\vspace{1em}

Dear Professors Alvelos, Carrasco Barra, Constantino, and Wei,

We are pleased to submit our manuscript ``The Cycle-Constrained Aerial
Suppression Base and Water Point Location Problem: Joint Siting under
Uncertainty with an Application to Cundinamarca, Colombia'' to the
International Transactions in Operational Research. As requested in
the call for papers, we indicate that this submission is intended for
the special issue on Operations Research in Wildfire Management.

The paper sits squarely within the call's scope, in particular the
listed topics of resource location-allocation in wildfire
preparedness, combined problems, budgeting, and suppression resource
operations, approached with stochastic programming and matheuristics.
It introduces, to our knowledge, the first model that sites aerial
suppression bases and water refill points jointly, with aircraft
productivity endogenous to the base to fire to water refill cycle, as
a two-stage mean-risk stochastic program based on conditional
value-at-risk (CVaR); a single budget covering bases, aircraft, and
water infrastructure makes the budget split endogenous. An exact
linearization and a fix-and-optimize matheuristic with provably exact
reduced subproblems make the full candidate catalog tractable. A
fully open-data, replicable case study for Cundinamarca, Colombia, is
built around the country's real Sikorsky S-70 Firehawk acquisition
and includes a sample average approximation replication analysis with
an optimality-gap estimate and an operational reading of the optimal
plan in real places, yielding policy-relevant structural findings on
when sequential planning suffices and on what limits containment.

The manuscript is original, unpublished, and not under consideration
elsewhere. The research received no external funding and the authors
declare no competing interests. All data sources are open and every
experiment is reproducible from the companion repository, available
to reviewers upon request.

Sincerely,

\vspace{1.5em}

Jaime Enrique Pesca Santos (corresponding author,
jaime.pesca@uexternado.edu.co)\\
Facultad de Administraci\'on de Empresas, Universidad Externado de
Colombia, Bogot\'a, Colombia

\vspace{0.8em}

Kiana Hikaru Ysa Morla (kiana.ysa@utec.edu.pe)\\
Universidad de Ingenier\'ia y Tecnolog\'ia (UTEC), Lima, Peru

\end{document}
